\subsubsection{开源代码逻辑、约束与选型边界}

代码流水线按下列接口拆分：
\begin{quote}
\small\ttfamily
geometry -> state -> rhs -> solver\\
-> validation -> export
\end{quote}
MATLAB 侧以 \texttt{ode45}、稀疏矩阵和优化工具箱形成对照实现；Python 侧
以 NumPy/SciPy 的数组、\texttt{solve\_ivp}、稀疏线性代数和网格数据结构实现。
每个模型提供一个可手算的小样例和一个真实规模样例，输出时间网格、状态量、
守恒残差、收敛曲线及下游所需的 CSV。

模型约束至少包括变量域、初边值、守恒关系、几何/运动学限制、参数物理范围和
数值稳定条件。能由规则直接计算的几何关系不使用神经网络；样本驱动关系没有
物理闭合时，不伪造微分方程。ODE 与 PDE 的边界取决于空间耦合是否不可忽略；
有限差分与有限元的边界取决于几何和介质复杂度；RK4 与自适应/刚性求解器的
边界取决于稳定域和尺度差异。

